% ===== CV TEMPLATE: PhD / Materials / Plasma Research =====
% Emphasis: experimental research credentials, thesis work, characterisation skills,
%           plasma-materials background. CTC role framed as research engineering.
% Claude Code instruction: copy to cv/cv.tex, then adjust section ordering or
% bullet emphasis to match the specific PhD project. Do not alter facts or dates.
\documentclass[a4paper,11pt]{article}
\usepackage[utf8]{inputenc}
\usepackage[T1]{fontenc}
\usepackage[margin=1in]{geometry}
\usepackage{setspace}
\usepackage{hyperref}
\usepackage{enumitem}
\setstretch{0.95}
\pagestyle{empty}

\begin{document}

%====================
% Header
%====================
\begin{center}
    {\LARGE \textbf{Bo Yuan}} \\
    \vspace{0.15cm}
    +46 769\,612\,787 \enspace $\cdot$ \enspace
    \href{mailto:boyua@kth.se}{boyua@kth.se} \enspace $\cdot$ \enspace
    Stockholm / Uppsala, Sweden
\end{center}

\vspace{0.3cm}

%====================
% Research Profile
%====================
\section*{Research Profile}

Experimental materials researcher with a background in microwave plasma systems and
plasma--material interactions. Currently conducting MSc thesis research at Uppsala
University on ion-irradiation-induced degradation of protective nitride coatings,
using a combination of ion-beam analysis, surface-sensitive diffraction, electron
microscopy, and mechanical characterisation. Prior to academia, led the full-cycle
R\&D and industrialisation of a high-power MPCVD system for diamond synthesis over
five years as CTO. Research interests focus on plasma-assisted materials processing,
ion--solid interactions, and structure--property relationships in functional materials.

%====================
% Education
%====================
\section*{Education}

\noindent\textbf{KTH Royal Institute of Technology}, Sweden \hfill 09/2024 -- 07/2026 (expected) \\
\textit{M.Sc.\ in Electromagnetics, Fusion and Space Engineering --- Plasma Track}

\smallskip
\noindent\textbf{Master's Thesis} \textit{(Ångström Laboratory / Tandem Laboratory, Uppsala University)} \\
\textit{Irradiation-induced degradation mechanisms in Cr--Nb--N coatings on Zr substrates}
\begin{itemize}[noitemsep, topsep=2pt, leftmargin=*]
    \item Independently designed ion irradiation experiments; calculated ion species,
          energies, and fluence using SRIM to reach target dpa levels
    \item Executed beamtime at the Tandem accelerator facility: sample preparation,
          mounting, and on-site monitoring
    \item Applied ToF-ERDA for depth-resolved compositional profiling; GI-XRD for
          near-surface structural characterisation; SEM/EDS for microstructural imaging;
          nanoindentation for hardness and elastic modulus measurements
    \item Carefully aligned SRIM damage depth profiles with analytical probe depths
          for quantitatively reliable interpretation of near-surface modifications
    \item Processed and interpreted all experimental datasets using Python
    \item Key finding: highest-content nanocrystalline/amorphous-phase CrNbN coating
          shows no significant degradation after irradiation and post-irradiation annealing
\end{itemize}

\smallskip
\noindent Relevant coursework: Plasma Physics, Electromagnetic Waves, Plasma--Surface
Interaction, Vacuum Technology, Material Selection and Characterisation,
Energy and Fusion Research, Microwave Modelling

\medskip
\noindent\textbf{Xi'an University of Technology}, China \hfill 09/2013 -- 07/2018 \\
\textit{B.Eng.\ in Electrical Engineering and Automation}

%====================
% Research Experience
%====================
\section*{Research Experience}

\noindent\textbf{Microwave Plasma CVD System Development} \hfill 2019 -- 2024 \\
\textit{Sichuan 330 Semiconductor Ltd.\ --- CTO / Lead Research Engineer}

\noindent Led the design, development, and experimental characterisation of a
high-power MPCVD system (2.45~GHz, 10~kW) for diamond synthesis, with a strong
focus on plasma physics and plasma--material interactions:

\begin{itemize}[noitemsep, topsep=2pt, leftmargin=*]
    \item Designed a TM$_{012}$-mode coaxial-tuned resonant cavity from first
          principles; optimised EM field distribution for uniform, stable plasma
          and enlarged effective deposition area to 75\,mm diameter
    \item Investigated the influence of microwave power density, cavity mode
          structure, gas chemistry (CH$_4$/H$_2$ ratios), and thermal conditions
          on plasma stability and carbon material growth behaviour
    \item Developed quantitative process parameter control enabling reproducible
          diamond growth: single-crystal up to 6\,mm thick on 7\,mm seed
          ($\geq$1 carat gem-quality); polycrystalline thin film on 60\,mm Si substrate
          (heatsink / optical window applications)
    \item Developed vacuum sealing solutions and leak detection procedures for
          high-vacuum plasma operation; designed and validated all vacuum-critical
          components
    \item Led transition from prototype to small-series production; delivered
          $\sim$20 complete systems; directed R\&D team (7 engineers)
\end{itemize}

\medskip
\noindent\textbf{Fusion Plasma Diagnostics --- Course Project (ED2246)} \hfill 2025 \\
\textit{KTH Royal Institute of Technology}
\begin{itemize}[noitemsep, topsep=2pt, leftmargin=*]
    \item Designed and executed plasma diagnostic experiments on a magnetic
          confinement fusion device at KTH
    \item Performed extensive experimental data processing and analysis of
          fusion plasma diagnostic signals
\end{itemize}

%====================
% Technical Skills
%====================
\section*{Technical Skills}

\noindent\textbf{Materials Characterisation:}
ToF-ERDA, GI-XRD, SEM with EDS, nanoindentation, SRIM simulations;
experience with accelerator-based facility workflows (Tandem Laboratory, Uppsala)

\smallskip
\noindent\textbf{Plasma and Experimental Systems:}
Microwave plasma systems (MPCVD, 2.45~GHz / 10~kW), plasma--surface interactions,
high-vacuum system design and operation, RF cavity design and tuning

\smallskip
\noindent\textbf{Simulation and Modelling:}
ANSYS HFSS ($\sim$2\,yr), COMSOL Multiphysics ($\sim$3\,yr), Keysight ADS ($\sim$1\,yr),
SRIM

\smallskip
\noindent\textbf{Data Analysis and Programming:}
Python (experimental data processing, analysis scripts, visualisation),
Matlab (basic); scientific data interpretation

\smallskip
\noindent\textbf{Hardware and CAD:}
SolidWorks ($\sim$10\,yr), Altium Designer ($\sim$10\,yr),
high-vacuum component design, Siemens S7-1200 PLC integration

%====================
% Patents
%====================
\section*{Patents}

\begin{itemize}[noitemsep, topsep=2pt, leftmargin=*]
    \item \textbf{Large-area compact microwave plasma CVD device}
          (CN212560433U, 2021) ---
          modular MPCVD architecture for stable large-area plasma generation
    \item \textbf{Waveguide sealing window for vacuum microwave transmission}
          (CN212392359U, 2021) ---
          vacuum-compatible dielectric window for high-power microwave systems
    \item \textbf{High-power TE--TEM microwave mode converter}
          (CN212162036U, 2020) ---
          low-loss mode conversion for efficient high-power microwave transmission
    \item \textbf{Microwave sliding short with choking structure}
          (CN212162044U, 2020) ---
          adjustable impedance tuning with reduced RF leakage
    \item \textbf{Tensioning device for MPCVD} (CN210906069U, 2020)
\end{itemize}

%====================
% Languages
%====================
\section*{Languages}

Chinese (native) \enspace $\cdot$ \enspace English (fluent, professional working level)

\end{document}
